\documentclass[manuscript,screen,review]{acmart}

%% Metadata
\acmJournal{TOMS}
\acmVolume{0}
\acmNumber{0}
\acmArticle{0}
\acmYear{2026}
\acmMonth{5}

\usepackage{amsmath,amssymb}
\usepackage{listings}
\usepackage{booktabs}
\usepackage{graphicx}
\usepackage{xcolor}

%% Code listings
\lstdefinelanguage{Rust}{
  keywords={fn, let, mut, pub, struct, impl, self, use, mod, match, if, else, for, while, return, true, false, cfg, feature},
  morecomment=[l]{//},
  morecomment=[s]{/*}{*/},
  morestring=[b]",
  sensitive=true,
}
\lstset{
  basicstyle=\ttfamily\small,
  keywordstyle=\color{blue}\bfseries,
  commentstyle=\color{gray}\itshape,
  stringstyle=\color{red},
  breaklines=true,
  frame=single,
  numbers=left,
  numberstyle=\tiny\color{gray},
}

\begin{document}

\title{rusty-SUNDIALS: A Memory-Safe Rust Implementation of CVODE with Formally Verified Newton Convergence}

\author{Xavier Callens}
\email{xavier.callens@example.com}
\affiliation{%
  \institution{Independent Researcher}
}

\begin{abstract}
We present \emph{rusty-SUNDIALS}, a from-scratch Rust implementation of the CVODE solver from the SUNDIALS suite (v7.4.0). Our implementation achieves behavioral equivalence with the LLNL C reference on the Robertson chemical kinetics benchmark while providing compile-time memory safety guarantees absent from the C original. Through a systematic, falsification-driven auto-research methodology, we identify and correct a critical convergence rate persistence bug in the nonlinear solver, achieving Newton iteration efficiency that \emph{exceeds} the C reference (1.40 iterations/step vs.\ 1.44). We formalize the C$\leftrightarrow$Rust behavioral equivalence via Lean~4 refinement proofs covering the nonlinear solver API surface. The solver is structured as a modular Rust workspace with zero \texttt{unsafe} blocks, builder-pattern configuration, and Cargo feature flags enabling experimental convergence heuristics. We discuss the architectural advantages of Rust's ownership model for ODE solver design and outline a path toward GPU acceleration on NVIDIA A100 hardware. All results are CI-validated on every commit across three platforms.
\end{abstract}

\begin{CCSXML}
<ccs2012>
<concept>
<concept_id>10002950.10003624</concept_id>
<concept_desc>Mathematics of computing~Ordinary differential equations</concept_desc>
<concept_significance>500</concept_significance>
</concept>
<concept>
<concept_id>10011007.10011074.10011099</concept_id>
<concept_desc>Software and its engineering~Software verification and validation</concept_desc>
<concept_significance>300</concept_significance>
</concept>
</ccs2012>
\end{CCSXML}

\ccsdesc[500]{Mathematics of computing~Ordinary differential equations}
\ccsdesc[300]{Software and its engineering~Software verification and validation}

\keywords{ODE solver, BDF methods, Rust, memory safety, SUNDIALS, Newton convergence, formal verification, Lean~4}

\maketitle

\section{Introduction}
The SUNDIALS suite~\cite{hindmarsh2005sundials} has been the \emph{de facto} standard for solving stiff and non-stiff ODEs for over two decades. Its CVODE component implements variable-order, variable-step BDF and Adams-Moulton methods, serving as the backbone of simulations in climate science, systems biology, and chemical engineering.

However, the C implementation carries inherent risks: null pointer dereferences, buffer overflows, and use-after-free bugs are structurally possible. Rust eliminates these defects at compile time while achieving C-equivalent performance.

\subsection{Contributions}
\begin{enumerate}
\item A complete Rust implementation of CVODE supporting BDF orders 1--5 and Adams-Moulton orders 1--12.
\item A falsification-driven auto-research methodology that tests convergence hypotheses against CI-validated benchmarks.
\item Discovery and correction of a convergence rate persistence bug causing 1.69$\times$ Newton iteration overhead.
\item Lean~4 formal verification of C$\leftrightarrow$Rust behavioral equivalence.
\item A GPU acceleration roadmap for NVIDIA A100 hardware.
\end{enumerate}

\section{The Convergence Rate Persistence Bug}

The most significant finding is a convergence rate persistence bug. LLNL's \texttt{cv\_crate} persists as a struct field; our V2 implementation reset it per-step.

\paragraph{Mathematical Impact.} The convergence test $\text{dcon} = \delta \cdot \min(1, \rho) / \text{tq}_4 \leq 1$ depends critically on~$\rho$. When $\rho$ persists ($\approx 0.01$): $\text{dcon} = \delta / 60$. When reset to 1.0: $\text{dcon} = \delta / 0.6$.

\begin{table}[h]
\caption{Impact of convergence rate persistence (CI-validated).}
\label{tab:results}
\begin{tabular}{lrrrr}
\toprule
\textbf{Metric} & \textbf{V2 (reset)} & \textbf{V3 (persistent)} & \textbf{C ref} & \textbf{V3/C} \\
\midrule
Steps           & 1,076 & 1,076 & 1,070 & 1.006$\times$ \\
Newton iters    & $\sim$2,603 & \textbf{1,503} & 1,537 & \textbf{0.98$\times$} \\
NI/step         & 2.42 & \textbf{1.40} & 1.44 & \textbf{0.97$\times$} \\
Conservation    & 8.88e-16 & 8.88e-16 & $\sim$1.1e-15 & Better \\
\bottomrule
\end{tabular}
\end{table}

\section{Benchmark Results}

\begin{table}[h]
\caption{Auto-research evolution across v11.x versions.}
\label{tab:evolution}
\begin{tabular}{llrrrr}
\toprule
\textbf{Version} & \textbf{Hypotheses} & \textbf{Steps} & \textbf{RHS} & \textbf{NI/step} & \textbf{Verdict} \\
\midrule
v11.1.0 & Baseline (FD Jac)    & 16,951 & 74,778 & ---  & Baseline \\
v11.2.0 & Analytical Jacobian  & 960    & 2,707  & ---  & Accept \\
v11.3.0 & H1--H3 Newton fixes  & 903    & 2,536  & ---  & Accept \\
v11.4.0 & H4--H6 corrected tq4 & 1,076  & 2,603  & 2.42 & Cond.\ Accept \\
v11.5.0 & H7--H8 persistent $\rho$ & 1,076 & 2,602 & \textbf{1.40} & \textbf{Accept} \\
\midrule
C ref   & LLNL SUNDIALS 7.4.0  & 1,070  & 1,537  & 1.44 & Reference \\
\bottomrule
\end{tabular}
\end{table}

\section*{Acknowledgments}
The auto-research methodology was developed with assistance from AI coding tools (Gemini, Antigravity by Google DeepMind). Peer reviews were conducted by Mistral AI (model: mistral-medium-latest). The use of generative AI is disclosed per ACM policy.

\bibliographystyle{ACM-Reference-Format}
\begin{thebibliography}{10}

\bibitem{hindmarsh2005sundials}
A.~C. Hindmarsh, P.~N. Brown, K.~E. Grant, S.~L. Lee, R.~Serban, D.~E. Shumaker, and C.~S. Woodward.
\newblock {SUNDIALS}: Suite of nonlinear and differential/algebraic equation solvers.
\newblock \emph{ACM Trans. Math. Softw.}, 31(3):363--396, 2005.

\bibitem{byrne1975polyalgorithm}
G.~D. Byrne and A.~C. Hindmarsh.
\newblock A polyalgorithm for the numerical solution of ordinary differential equations.
\newblock \emph{ACM Trans. Math. Softw.}, 1(1):71--96, 1975.

\bibitem{matsakis2014rust}
N.~D. Matsakis and F.~S. Klock~II.
\newblock The {R}ust language.
\newblock \emph{ACM SIGAda Ada Letters}, 34(3):103--104, 2014.

\bibitem{demoura2021lean4}
L.~de~Moura and S.~Ullrich.
\newblock The {L}ean~4 theorem prover and programming language.
\newblock In \emph{CADE-28}, 2021.

\bibitem{robertson1966solution}
H.~H. Robertson.
\newblock The solution of a set of reaction rate equations.
\newblock In J.~Walsh, editor, \emph{Numerical Analysis: An Introduction}, pages 178--182. Academic Press, 1966.

\end{thebibliography}

\end{document}
